% ============================================================
%  ANEXOS
% ============================================================

% -----------------------------------------------------------
\subsection*{Anexo A: Acceso a la aplicación web}
\label{anexo:aplicacion}
\addcontentsline{toc}{subsection}{Anexo A: Acceso a la aplicación web}
% -----------------------------------------------------------

El sistema CAIQ está desplegado como aplicación web de acceso público en Hugging Face
Spaces. La aplicación puede consultarse en la siguiente URL:

\begin{center}
  \url{https://huggingface.co/spaces/relan02/caiq}
\end{center}

No se requiere registro ni autenticación. El único requisito es disponer de un currículum
vitae en formato PDF, DOCX o TXT. El corpus de ofertas de empleo subyacente se actualiza
automáticamente cada semana mediante el \textit{pipeline} descrito en la
Sección~\ref{subsec:pipeline_scraping}.

La aplicación se organiza en tres vistas principales accesibles desde el menú lateral:
la vista de \textbf{resumen del perfil} (competencias detectadas, nivel de
\textit{seniority} y distribución del \textit{skill gap} por banda de prioridad), la
vista de \textbf{recomendaciones} (cursos, másteres y ofertas ordenados por su
contribución al cierre de la brecha) y la vista de \textbf{inteligencia de
\textit{skills}} (demanda de competencias en el mercado filtrable por familia de rol).
Las capturas de estas tres vistas se presentan en las Figuras~13--15 del
capítulo de Resultados.

% -----------------------------------------------------------
\subsection*{Anexo B: Estructura del código fuente}
\label{anexo:codigo}
\addcontentsline{toc}{subsection}{Anexo B: Estructura del código fuente}
% -----------------------------------------------------------

El código fuente del sistema se organiza en cuatro módulos que corresponden a las capas
funcionales descritas en la Sección~\ref{subsec:arquitectura_general_del_sistema}.

\textbf{Pipeline de \textit{scraping} e ingesta} (\texttt{actualizar\_ofertas.py}).
Script principal de actualización del corpus. Implementa el flujo completo: scraping
multi-fuente mediante \texttt{JobSpy}, deduplicación por URL, fusión incremental con el
histórico, caducidad temporal configurable y verificación activa de URLs con caché
persistente (\texttt{job\_url\_status.csv}). Se invoca semanalmente de forma automática
mediante la tarea programada definida en \texttt{setup\_tarea\_programada.ps1} (todos
los lunes a las 08:00).

\textbf{Pipeline ETL y capa semántica} (\texttt{01\_etl/pipelines/}). Conjunto de
scripts numerados que implementan el proceso de transformación y enriquecimiento:
limpieza y normalización (\texttt{01\_etl\_curated\_layer.py}), generación de
\textit{embeddings} y tabla \texttt{role\_skill\_demand}
(\texttt{02\_build\_semantic\_features.py}), entrenamiento y calibración del motor de
recomendación (\texttt{03\_train\_and\_tune\_recommender.py}), evaluación cuantitativa
(\texttt{04\_evaluate\_recommender\_metrics.py}) y estudio de ablación de variantes
(\texttt{05\_ablation\_study\_recommender.py}).

\textbf{Aplicación web} (\texttt{app/caiq/}). Contiene el código de la
interfaz \texttt{Streamlit}: \texttt{caiq\_app.py} (módulo principal),
\texttt{config/skills\_taxonomy.json} (taxonomía controlada de 71 competencias con 477
alias léxicos), \texttt{Dockerfile} (imagen de despliegue) y \texttt{requirements.txt}
(dependencias de producción).

\textbf{Validación} (\texttt{docs/validation/}). Scripts de validación cuantitativa sobre
el corpus de 500 CVs: extracción de competencias, comparación contra anotación manual
($n$\,=\,87) y generación de las métricas de precisión, \textit{recall} y F1 reportadas
en la Sección~\ref{subsec:resultados_obj4}.

% -----------------------------------------------------------
\subsection*{Anexo C: Taxonomía de competencias}
\label{anexo:taxonomia}
\addcontentsline{toc}{subsection}{Anexo C: Taxonomía de competencias}
% -----------------------------------------------------------

La Tabla~\ref{tab:taxonomia_completa} recoge las 71 competencias que componen la
taxonomía controlada del sistema, con el número de alias léxicos definidos para cada una.
Los alias permiten reconocer la misma competencia bajo las distintas denominaciones
habituales en el mercado laboral (por ejemplo, \texttt{sql} incluye variantes como
\textit{MySQL}, \textit{PostgreSQL}, \textit{T-SQL} o \textit{PL/SQL}).

\begin{table}[H]
  \caption{Taxonomía de competencias del sistema CAIQ: 71 competencias canónicas y
  número de alias léxicos.}
  \label{tab:taxonomia_completa}
  \centering
  \footnotesize
  \begin{tabular}{|p{4.0cm}|c||p{4.0cm}|c|}
    \hline
    \rowcolor{gray!15}
    \textbf{Competencia} & \textbf{Alias} & \textbf{Competencia} & \textbf{Alias} \\
    \hline
    a/b testing         &  3 & lightgbm           &  3 \\
    airflow             &  5 & linux              &  9 \\
    api                 &  8 & llm                & 10 \\
    aws                 & 11 & looker             &  4 \\
    azure               &  9 & machine learning   & 17 \\
    big data            &  3 & mapreduce          &  3 \\
    business intelligence &  8 & mathematics      & 11 \\
    c++                 &  3 & microsoft fabric   &  4 \\
    ci/cd               &  9 & mlflow             &  6 \\
    computer vision     & 10 & mlops              & 11 \\
    dask                &  4 & nlp                & 14 \\
    data analysis       &  4 & nosql              &  9 \\
    data engineering    &  3 & numpy              &  4 \\
    data science        &  3 & pandas             &  3 \\
    data visualization  & 16 & power bi           &  4 \\
    data governance     &  8 & python             &  4 \\
    databricks          &  5 & pytorch            &  3 \\
    dbt                 &  4 & r                  & 14 \\
    deep learning       & 17 & redshift           &  3 \\
    docker              &  7 & sas                &  5 \\
    elasticsearch       &  7 & scala              &  4 \\
    etl                 &  6 & scikit-learn       &  5 \\
    excel               &  8 & snowflake          &  3 \\
    fastapi             &  8 & spark              &  6 \\
    feature engineering &  5 & spss               &  4 \\
    forecasting         &  4 & sql                & 18 \\
    gcp                 & 10 & statistics         & 16 \\
    generative ai       & 17 & tableau            &  6 \\
    git                 &  7 & talend             &  3 \\
    hadoop              &  3 & tensorflow         &  4 \\
    hdfs                &  3 & terraform          &  6 \\
    hive                &  4 & time series        &  9 \\
    informatica         &  4 & xgboost            &  3 \\
    java                &  4 & yarn               &  3 \\
    kafka               &  7 & keras              &  4 \\
    kubernetes          &  5 & & \\
    \hline
    \multicolumn{4}{|r|}{\textbf{Total: 71 competencias — 477 alias}} \\
    \hline
  \end{tabular}
\end{table}
